\documentclass[aspectratio=169,10pt]{beamer}
\usetheme[numbering=fraction,progressbar=frametitle]{metropolis}
\usepackage{fontspec}
\setsansfont{Latin Modern Sans}
\usepackage{amsmath,amssymb,mathtools,graphicx,booktabs}
\usepackage{polyglossia}
\setdefaultlanguage{english}
\setbeamertemplate{navigation symbols}{}
\setbeamerfont{frametitle}{size=\large}
\setbeamerfont{normal text}{size=\normalsize}
\setlength{\parskip}{.45em}
\hypersetup{pdfauthor={},pdfsubject={Lecture slides},pdfkeywords={},pdftitle={GCD, Bézout's and Gauss's theorems}}
\providecommand{\N}{\mathbb{N}}
\providecommand{\Z}{\mathbb{Z}}
\providecommand{\D}{\mathbb{D}}
\providecommand{\Q}{\mathbb{Q}}
\providecommand{\R}{\mathbb{R}}
\providecommand{\C}{\mathbb{C}}
\providecommand{\vect}[1]{\overrightarrow{#1}}
\providecommand{\vi}{\vec{\imath}}
\providecommand{\vj}{\vec{\jmath}}
\providecommand{\vk}{\vec{k}}
\providecommand{\norm}[1]{\left\lVert #1\right\rVert}
\providecommand{\vc}[2]{\begin{pmatrix}#1\\#2\end{pmatrix}}
\providecommand{\deter}[2]{\det\left(#1,#2\right)}
\providecommand{\abs}[1]{\left\lvert #1\right\rvert}
\providecommand{\dd}{\,\mathrm{d}}
\providecommand{\e}{\mathrm{e}}
\providecommand{\Card}{\operatorname{Card}}
\providecommand{\E}{\operatorname{E}}
\providecommand{\V}{\operatorname{V}}
\providecommand{\degre}{^{\circ}}
\providecommand{\un}[1]{\,\mathrm{#1}}
\providecommand{\sci}[2]{#1\times10^{#2}}
\providecommand{\pgcd}{\operatorname{PGCD}}
\providecommand{\ppcm}{\operatorname{PPCM}}
\title{GCD, Bézout's and Gauss's theorems}
\subtitle{Lecture and worked examples}
\date{}
\author{}
\begin{document}
\begin{frame}[plain]
\vfill
{\usebeamerfont{title}\inserttitle\par}
\medskip
{\usebeamerfont{subtitle}\insertsubtitle\par}
\medskip
{\usebeamercolor[fg]{progress bar}\rule{\textwidth}{.5pt}}
\vfill
\end{frame}
\begin{frame}{Objectives}
Find a gcd with Euclid's algorithm; obtain a Bézout pair; use Bézout's and Gauss's theorems; find an inverse modulo \(n\); solve an equation \(ax+by=c\) in integers.
\end{frame}
\begin{frame}{Introductory activity}
We have a \(9\)-litre can, a \(7\)-litre can, both without markings, and a water supply. We may fill a can, empty it, or pour one into the other until one is empty or the other is full.
\end{frame}
\begin{frame}{Introductory activity (2)}
\begin{enumerate}
\setcounter{enumi}{0}\item Show that one can get exactly \(2\) litres in the large can.

\item If we pour in \(u\) times \(9\) litres altogether and throw away \(v\) times \(7\) litres, the quantity left is \(9u-7v\). Find natural numbers \(u\) and \(v\) such that \(9u-7v=1\). Deduce how to get \(1\) litre.

\item The cans are replaced by cans of \(6\) and \(9\) litres. Show that every quantity of the form \(6u+9v\), with \(u\) and \(v\) integers, is a multiple of \(3\). Can one get \(1\) litre?

\end{enumerate}
\end{frame}
\begin{frame}{Introductory activity (3)}
\begin{enumerate}
\setcounter{enumi}{3}\item With cans of \(a\) and \(b\) litres, which whole quantities does it seem possible to get? State a conjecture involving the greatest common divisor of \(a\) and \(b\).

\end{enumerate}
\end{frame}
\begin{frame}{GCD and Euclid's algorithm}
Let \(a\) and \(b\) be two integers, not both zero. The set of their positive common divisors is finite and contains \(1\); its largest element is called the \emph{greatest common divisor} of \(a\) and \(b\), written \(\operatorname{gcd}(a,b)\). We say that \(a\) and \(b\) are \emph{coprime} if \(\operatorname{gcd}(a,b)=1\).
\end{frame}
\begin{frame}{GCD and Euclid's algorithm}
Let \(a\) and \(b\) be two integers with \(b\neq0\). For every integer \(k\), the common divisors of \(a\) and \(b\) are those of \(b\) and \(a-kb\). In particular, if \(r\) is the remainder of the division of \(a\) by \(b\ge1\), then \(\operatorname{gcd}(a,b)=\operatorname{gcd}(b,r)\).
\end{frame}
\begin{frame}{GCD and Euclid's algorithm}
If \(d\mid a\) and \(d\mid b\), then \(d\mid a-kb\). Conversely, if \(d\mid b\) and \(d\mid a-kb\), then \(d\mid(a-kb)+kb=a\). The two sets of common divisors are equal, so their largest elements are too.
\end{frame}
\begin{frame}{GCD and Euclid's algorithm}
Replace \((a,b)\) by \((b,r)\) as long as the remainder \(r\) is not zero. The remainders form a strictly decreasing sequence of natural numbers, so the algorithm stops. The gcd is the last nonzero remainder. Example:
\[1{,}053=3\times325+78,\qquad325=4\times78+13,\qquad78=6\times13+0,\]
so \(\operatorname{gcd}(1{,}053,325)=13\).
\end{frame}
\begin{frame}{Bézout's theorem}
Let \(a\) and \(b\) be two integers, not both zero. There are integers \(u\) and \(v\) such that \(au+bv=\operatorname{gcd}(a,b)\). Moreover, every common divisor of \(a\) and \(b\) divides \(\operatorname{gcd}(a,b)\).
\end{frame}
\begin{frame}{Bézout's theorem}
Let \(E\) be the set of positive integers of the form \(au+bv\), with \(u\), \(v\) integers. \(E\) is not empty (it contains \(a^2+b^2\)). Let \(d=au_0+bv_0\) be its smallest element.

Divide \(a\) by \(d\): \(a=dq+r\) with \(0\le r<d\). Then \(r=a-q(au_0+bv_0)=a(1-qu_0)+b(-qv_0)\) is of the form \(au+bv\). If it were positive, it would belong to \(E\) while being smaller than \(d\): impossible. So \(r=0\) and \(d\mid a\). Likewise \(d\mid b\).
\end{frame}
\begin{frame}{Bézout's theorem (2)}
Finally every common divisor \(c\) of \(a\) and \(b\) divides \(au_0+bv_0=d\), so \(c\le d\). Hence \(d\) is the greatest common divisor: \(d=\operatorname{gcd}(a,b)\).
\end{frame}
\begin{frame}{Bézout's theorem}
Two integers \(a\) and \(b\) are coprime if and only if there are integers \(u\) and \(v\) such that \(au+bv=1\).
\end{frame}
\begin{frame}{Bézout's theorem}
If \(\operatorname{gcd}(a,b)=1\), the previous theorem provides \(u\) and \(v\). Conversely, if \(au+bv=1\), every positive common divisor of \(a\) and \(b\) divides \(1\), hence equals \(1\).
\end{frame}
\begin{frame}{Bézout's theorem}
With the previous example: \(13=325-4\times78=325-4(1{,}053-3\times325)=13\times325-4\times1{,}053\). A suitable pair for \((1{,}053,325)\) is \((-4,13)\).
\end{frame}
\begin{frame}{Bézout's theorem}
The following function returns \((d,u,v)\) with \(d=\operatorname{gcd}(a,b)=au+bv\), for natural numbers \(a\) and \(b\) not both zero.
lstlisting
def euclide\_etendu(a, b):
    u0, v0, u1, v1 = 1, 0, 0, 1
    while b != 0:
        q, r = a // b, a \% b
        a, b = b, r
        u0, u1 = u1, u0 - q * u1
        v0, v1 = v1, v0 - q * v1
    return a, u0, v0
lstlisting
At every step, the two current remainders can be written \(au_0+bv_0\) and \(au_1+bv_1\) in terms of the initial values of \(a\) and \(b\).
\end{frame}
\begin{frame}{Gauss's theorem}
Let \(a\), \(b\), \(c\) be three integers. If \(a\mid bc\) and \(a\) and \(b\) are coprime, then \(a\mid c\).
\end{frame}
\begin{frame}{Gauss's theorem}
By Bézout's theorem, there are \(u\) and \(v\) such that \(au+bv=1\). Multiplying by \(c\): \(c=acu+bcv\). Now \(a\mid acu\) and \(a\mid bc\), so \(a\) divides \(acu+bcv=c\).
\end{frame}
\begin{frame}{Gauss's theorem}
If \(a\) and \(b\) are coprime and both divide \(n\), then \(ab\) divides \(n\).
\end{frame}
\begin{frame}{Gauss's theorem}
\(n=ak\). Since \(b\mid ak\) and \(\operatorname{gcd}(a,b)=1\), Gauss gives \(b\mid k\): \(k=bk'\) and \(n=abk'\).
\end{frame}
\begin{frame}{Gauss's theorem}
Let \(n\ge2\) and \(a\) an integer. There is an integer \(u\) such that \(au\equiv1\pmod n\) if and only if \(\operatorname{gcd}(a,n)=1\). Such a \(u\) is an inverse of \(a\) modulo \(n\); it allows us to solve \(ax\equiv b\pmod n\): \(x\equiv ub\pmod n\).
\end{frame}
\begin{frame}{Gauss's theorem}
\(au\equiv1\pmod n\) is equivalent to the existence of an integer \(v\) such that \(au+nv=1\): this is Bézout's theorem.
\end{frame}
\begin{frame}{Gauss's theorem}
\begin{enumerate}
\setcounter{enumi}{0}
\item It has solutions if and only if \(\operatorname{gcd}(a,b)\) divides \(c\); we then divide by this gcd.

\item Find a particular solution \((x_0,y_0)\).

\item By subtraction, \(a(x-x_0)=-b(y-y_0)\); with \(a\) and \(b\) coprime, Gauss gives \(x-x_0=bk\), then \(y-y_0=-ak\).

\item Check that the pairs \((x_0+bk,y_0-ak)\), \(k\in\Z\), are indeed solutions.

\end{enumerate}

Example: \(7x+5y=1\) has the particular solution \((-2,3)\); its solutions are the pairs \((-2+5k,3-7k)\), \(k\in\Z\).
\end{frame}
\end{document}
